\section{Operational use case views}
The customer view is complemented by Figures~\ref{fig:operations} and~\ref{fig:backoffice}. Catalogue work belongs to the Merchandiser; receiving, packing, and physical inspection belong to Warehouse Staff. Customer Service participates in return decisions, while the Shipping Provider supports shipment registration and supplies delivery events. Finance Staff reconcile refunds and financial reports. The Manager reads business reports; the Administrator manages access rather than being granted every operational responsibility automatically. UC17 notifications and UC18 optional reviews are supporting services specified in the catalogue; they do not change the stock and money commitments shown in the core views.
\landdiagram{02-operations-use-cases}{Catalogue, warehouse, and delivery responsibilities.}{fig:operations}
\diagram{16-backoffice-use-cases}{Customer service, finance, management, and access administration.}{fig:backoffice}
Figure~\ref{fig:support-uc} makes the supporting goals explicit in the native model. A Customer receives committed transaction notifications and may submit an eligible purchased-product review. Finance Staff can reconcile unresolved payment events through UC05; the provider remains the normal external event source. Review publication is optional under the Could priority.
\diagram{18-support-use-cases}{Transaction notices, optional reviews, and manual financial reconciliation.}{fig:support-uc}

\section{Supporting activity and swimlane models}
Figure~\ref{fig:receiving} develops the stock receipt obligation. A receipt key prevents a repeated submission from becoming another stock increase. The actual count separates damaged from saleable units. Only validated new data reaches the atomic posting action; the alternative reports invalid input or the already recorded result.
\diagram{14-inventory-receiving}{Receiving saleable stock with validation and duplicate protection.}{fig:receiving}

Figures~\ref{fig:inspection} and~\ref{fig:remedy} divide post-purchase service into physical verification and commercial resolution. The inspection workflow begins with an already approved request. Warehouse Staff verify the parcel and condition, then the system checks quantities and receipt identity before committing stock changes. A refund decision alone cannot produce this physical receipt evidence. A disputed receipt ends the current processing attempt without committing stock; staff must resolve the discrepancy before another attempt.
\diagram{05-return-inspection}{Return inspection and controlled restoration of saleable stock.}{fig:inspection}

Once goods have been accepted and inspected, an exchange preference is resolved against actual replacement stock. Figure~\ref{fig:remedy} reserves the requested equal-price replacement before confirming linked fulfilment. Unavailable stock produces a pending customer decision; the final node terminates that processing attempt, not the entire ReturnRequest lifecycle. A later refund choice resumes UC14, while withdrawal closes the exchange preference. The refund branch schedules financial work and does not assert that the provider has already completed it.
\diagram{15-return-remedy}{Commercial resolution after inspection: refund or reserved exchange.}{fig:remedy}

\section{Order lifecycle and transaction states}
Figure~\ref{fig:order-state} describes \emph{order fulfilment}, not every financial or post-purchase fact. The initial transition enters PendingPayment only after checkout has committed the order and its reservations. Confirmed means trusted capture and valid allocation are both established. Packing blocks customer cancellation in the first-release policy. Dispatch consumes stock allocation; verified delivery records the delivery fact. Cancelled, Expired, Delivered, and DeliveryFailed are persistent business states. They do not need to destroy the order object or erase evidence.
\diagram{06-order-state}{Order fulfilment lifecycle; money and returns have separate states.}{fig:order-state}
\begin{longtable}{P{.19\linewidth}P{.28\linewidth}P{.42\linewidth}}
\caption{Order transition conditions and effects}\label{tab:order-transitions}\\
\toprule Event & Allowed source and guard & Result and invariant \\\midrule\endfirsthead
\toprule Event & Allowed source and guard & Result and invariant \\\midrule\endhead
verifiedCapture & PendingPayment; matching trusted capture; hold still active & Enter Confirmed; retain allocations for fulfilment and remove the initial payment expiry condition.\\
startPacking & Confirmed; verified payment and sufficient active allocation & Enter Packing under a lock/version check; a competing cancellation cannot also win.\\
dispatch & Packing; SKU checks complete; shipment reference confirmed & Enter Shipped; decrement saleable on-hand and reserved by the dispatched quantities once.\\
verifiedDelivery & Shipped; validated carrier evidence & Enter Delivered; store delivery time used by the return policy.\\
closeFailure & Shipped; Customer Service authorises permanent non-delivery using verified carrier evidence & Enter DeliveryFailed; schedule captured-funds refund independently. Consumed stock is not automatically restored.\\
cancel & PendingPayment or Confirmed; packing has not begun & Enter Cancelled; release active stock allocation; separately cancel unresolved payment or request refund of captured funds.\\
holdExpired & PendingPayment; initial 15-minute deadline passed & Enter Expired; release active holds once; reconcile provider status and compensate any late capture.\\
\bottomrule
\end{longtable}

Payment status uses Pending, Processing, Succeeded, Failed, and Cancelled. Failed may be retried while the same order remains eligible; unresolved Processing is reconciled instead of treated as failure. A provider object that already succeeded cannot be moved back to Failed by an older notification. Captured money remains recorded even when fulfilment is Cancelled or Expired.

Shipment status uses Registered, InTransit, Delivered, DeliveryFailed, and ReturnedToSender. A temporary failed attempt remains InTransit while redelivery is possible. Permanent non-delivery requires verified evidence and an authorised UC12 closure; it cannot overwrite Delivered. The first-release policy refunds the full captured order amount, including delivery, for this closure. A parcel physically returned by the carrier is counted and inspected through UC10 using a unique shipment/receipt reference; only saleable units are posted back. Lost parcels produce no stock increase. This is separate from a customer's UC08 return of an already delivered order.

Refund status uses Requested, Pending, Succeeded, and Failed. Requested/Pending amounts reserve a financial entitlement. Failed means an exception needing verified resolution; retry of that refund keeps the same commitment and request key. A commitment is released only when cancellation/failure is authoritative and the business decision no longer requires the refund. Returns use Requested, Approved, Received, Inspected, AwaitingChoice, Resolved, Rejected, and Withdrawn. Split receipts may accumulate while the request remains Received; Resolved requires the chosen commercial remedy to finish. These states are independent of the original order remaining Delivered.

\section{Domain concepts and analysis class models}
The class models identify business concepts and their relationships. Identifier denotes a stable business identity; MoneyVND denotes an exact integer number of VND; DateTime denotes an instant stored consistently. They are conceptual types rather than claims about a programming language. Attributes are representative business information; the database design expands them with keys, audit, and integration fields.

Figure~\ref{fig:catalogue-class} establishes Product as a garment design and ProductVariant as a sellable SKU. Each variant belongs to one Product and has one InventoryBalance in the one-warehouse case. Every Product belongs to one primary Category; a category can contain many products. A SizeChart can be reused by several products, while a product has at most one applicable chart version. Accessories may have no size chart. The zero-or-more variant multiplicity permits draft products; publication separately requires at least one active, valid variant. InventoryBalance.available is a derived view of onHand minus reserved, not an independently editable stock counter.
\landdiagram{07-catalogue-class}{Catalogue, size charts, SKU variants, and warehouse inventory.}{fig:catalogue-class}

Figure~\ref{fig:order-class} separates cart intention from purchase commitment. A Customer has at most one active Cart and may have many Orders. Each CartLine selects one exact ProductVariant. An Order has at least one OrderLine and exactly one AddressSnapshot, while every OrderLine keeps a link to the retained variant plus an immutable commercial description. Product edits cannot silently rewrite the purchased title, size, colour, chart version, price, discount, or delivery details.
\landdiagram{08-order-class}{Customers, carts, committed order lines, and address snapshots.}{fig:order-class}

Figure~\ref{fig:reservation-class} makes stock commitment explicit. A Reservation belongs to one original Order and one SKU. Its kind distinguishes initial purchase allocation from a linked replacement allocation; a replacement additionally references its ExchangeFulfilment in persistence. Active purchase holds expire while PendingPayment; after confirmation they remain allocated until consumption or permitted cancellation. A replacement does not change the original order's Delivered state. Reservation histories are retained after release/consumption, so a historical order may have several reservations even though only its active allocations contribute to reserved stock.
\landdiagram{17-reservation-class}{Stock reservations are explicit commitments and are independent of carts.}{fig:reservation-class}

Figure~\ref{fig:aftersales-class} separates payment attempts, shipment evidence, returned quantities, and refunds. An Order can have several payment attempts, but only one successful payment object is accepted for its original payable total; further real captures require compensation. Each Refund references one Payment. An original Order has at most one Shipment in the first release. Each ReturnLine identifies one original OrderLine; a ReturnRequest groups one or more such lines. A ReturnLine may lead to one linked ExchangeFulfilment for its allocated replacement quantity. Additional persistence links identify replacement SKU, tracking, reservations, and refund allocation; they do not turn the analysis model into a diagram of every database column.
\landdiagram{09-aftersales-class}{Independent payment, shipment, return, and exchange concepts.}{fig:aftersales-class}

\section{Cross-model consistency}
The shared ProductVariant, Order, and OrderLine concepts carry the same meaning in every view. UC04 and the checkout sequence both establish stock commitments before remote payment processing. The order state machine permits packing only from Confirmed; UC11 performs that guarded transition. The inspection swimlane restores stock from physical evidence, whereas UC14 changes financial state. UC08 and the class model support partial per-line returns; the quantity rules restrict the otherwise broad many-valued associations. Multiplicity alone cannot express monetary caps, cumulative return quantities, stock inequalities, or time windows, so these remain explicit business rules and transaction constraints.
